\section{Canonical encodings and wire format}
\label{app:wire}

\subsection{Arithmetic hash domains}

The Poseidon2 sponge domains for owner derivation, nullifiers, and notes are,
respectively,
\[
 \mathtt{0x41530001},\qquad
 \mathtt{0x41530002},\qquad
 \mathtt{0x41530003}
 \quad\text{in }\mathbb F.
\]
The atomic-v3 Merkle compressor uses the width-16 state
\(L\|R\), adds \(\mathtt{0x41531005}\) to lane 15, applies the fixed
Poseidon2 permutation, and returns lanes 0 through 7.  The statement digest is
\[
 \operatorname{SHA256}(
   \mathtt{aspis/atomic\mbox{-}payment\mbox{-}statement/v3}\|\mathsf{payload}).
\]

\subsection{Public statement payload}

\begin{table}[ht]
\centering
\caption{Canonical atomic-v3 statement payload.}
\label{tab:statement-wire-appendix}
\begin{tabular}{@{}rrll@{}}
\toprule
offset & bytes & field & encoding \\
\midrule
0 & 1 & version & \(3\) \\
1 & 1 & tree depth & \(20\) \\
2 & 6 & reserved & all zero \\
8 & 32 & pool address \(P\) & raw public key \\
40 & 8 & sequence \(s\) & little-endian \(u64\) \\
48 & 32 & current anchor \(A\) & eight canonical M31 limbs \\
80 & 32 & nullifier \(\nu\) & eight canonical M31 limbs \\
112 & 32 & output commitment \(C_{\rm out}\) & eight canonical M31 limbs \\
144 & 32 & output anchor \(A'\) & eight canonical M31 limbs \\
176 & 4 & asset \(a\) & canonical little-endian M31 \\
180 & 4 & fee \(f\) & little-endian \(u32<2^{30}\) \\
\bottomrule
\end{tabular}
\end{table}

Changing any field changes the statement digest.  The six reserved bytes,
version, and depth are checked rather than ignored.

\subsection{Proof prefix}

The fixed portion of a \Profile{} proof has 6785 bytes.  Its physical order
is listed in Table~\ref{tab:proof-prefix-wire}.  The transcript absorbs fields
in the logical order of Section~\ref{sec:construction}; in particular, the
batch and fold work records are logically positioned before dependent
challenges even though the three work arrays occupy a compact suffix.

\begin{table}[ht]
\centering
\small
\caption{\Profile{} fixed proof prefix.  End offsets are exclusive.}
\label{tab:proof-prefix-wire}
\begin{tabular}{@{}rrrl@{}}
\toprule
start & end & bytes & object \\
\midrule
0 & 16 & 16 & canonical v4-s2 header \\
16 & 48 & 32 & public mask nonce \\
48 & 80 & 32 & \(C_1\) root \\
80 & 112 & 32 & \(C_2\) root \\
112 & 128 & 16 & initial masked claim \\
128 & 4608 & 4480 & ten-round masked sumcheck \\
4608 & 5952 & 1344 & 84 statement-point evaluations \\
5952 & 6096 & 144 & round-0 OOD values and relation sumcheck \\
6096 & 6272 & 176 & round-1 root, OOD values, and sumcheck \\
6272 & 6448 & 176 & round-2 root, OOD values, and sumcheck \\
6448 & 6624 & 176 & round-3 root, OOD values, and sumcheck \\
6624 & 6688 & 64 & four final-polynomial coefficients \\
6688 & 6696 & 8 & final-work nonce \\
6696 & 6728 & 32 & four fold-work nonces \\
6728 & 6736 & 8 & batch-work nonce \\
6736 & 6784 & 48 & three \(D\)-lane claims \\
6784 & 6785 & 1 & query selector in \(\{0,1,2\}\) \\
\bottomrule
\end{tabular}
\end{table}

The header fixes profile 23, \(\log_2\) message size 10, blowup logarithm 9,
18 queries, four rounds, raw-fiber fold payloads, radix-four minimal-subtree
Merkle mode, final polynomial logarithmic length 2, and the exact flag set.
All field encodings in the prefix are canonical.

\subsection{Private opening sections}

The suffix consists of five sections in the order \(C_1,C_2,W_1,W_2,W_3\).
Their Merkle depths and fixed opened-value widths are
\[
 (17,17,15,13,11),\qquad
 (416,192,64,64,64)\ \text{bytes}.
\]
Each section is encoded as
\[
 \mathtt{count}_{u16}\ \|
 (\mathtt{value}_{\rm fixed}\|\mathtt{salt}_{32})^{\mathtt{count}}\ \|
 \mathtt{frontier\_count}_{u32}\ \|
 \mathtt{frontier}_{32}^{\mathtt{frontier\_count}}.
\]
The transcript-derived indices are not serialized.  The parser derives the
layer-zero indices and their successive right shifts, requires their
canonical order, authenticates every value--salt record and frontier, and
rejects trailing bytes.  Consequently proof length varies with the minimal
frontier induced by the public schedule.

\subsection{Proof-account frame}

The exact account image before finalization is
\[
 \mathtt{ASPU}\|\mathtt{proof\_len}_{u32,\rm le}
 \|\mathtt{upload\_authority}_{32}\|\mathtt{proof}.
\]
Authorized chunk writes affect only the declared proof payload.  The
\textsf{FinalizeProof} instruction (wire tag 62) replaces bytes 8 through 39
with zero and preserves the payload.  Production verification accepts only
the all-zero sentinel.  The atomic pool is initialized separately by wire tag
63.  Wire tag 60 retains the proof account; the evaluated mainnet transition
uses wire tag 65 to verify, mutate state, and refund the proof account
atomically.  Wire tag 64 closes a separately finalized proof account without
applying a spend.

\section{Transcript boundary and imported-theorem checks}
\label{app:transcript-hypotheses}

The soundness evaluator treats the transcript as 32 prover-message
boundaries for the BCS transform.  These boundaries are fixed by the parser;
changing the prefix shape, sumcheck round count, OOD records, commitment
rounds, or selector position invalidates the parameter fingerprint rather
than silently changing \(R\).

The imported list-decoding/MCA step is instantiated only after the following
local checks.
\begin{enumerate}
  \item The layer-zero domain contains \(2^{19}\) circle symbols and its
        \(2^{17}\) arity-four fibers map injectively to line roots under
        \(t=2x^2-1\); no root is one.
  \item The message has \(2^{10}\) coefficients, giving rate \(1/512\), and
        the selected proximity threshold lies in the proven Johnson regime.
  \item The width-29 \(\gamma\)-batch is a polynomial-generator family over
        the fixed extension field, with \(\gamma\ne0\).
  \item The two commitment phases and four arity-four folds use the same
        coordinate scaling as the circle-to-GRS transport.
  \item Queries are an ordered sample of 18 distinct fibers; the verifier and
        numerical ledger use without-replacement probabilities.
  \item Two OOD samples and the relation sumcheck are present at every fold
        layer, and the final object has exactly four coefficients.
\end{enumerate}
These checks supply the hypotheses of the polynomial-generator MCA and circle
transport results used in
Lemma~\ref{lem:johnson-mca-batch}~\cite{BordageChiesaGuanManzur2025,CarmonEtAl2026Stwo}.

\section{Complete soundness arithmetic}
\label{app:soundness-ledger}

The four individual fold-event bounds are
\[
 109.52994269233842,\quad
 111.22628180453175,\quad
 112.85813508022261,\quad
 113.84064801383485
\]
bits.  Their union is the 108.98543226575069-bit row in
Table~\ref{tab:soundness-events}.  For one query branch, replacing the
selector-local final-miss row by 111.76870163436465 bits and unioning all
sixteen rows gives
\[
 -\log_2\varepsilon_{\rm round}^{(0)}
 =106.79080600295417.
\]
The endpoint reconstruction is
\begin{align*}
 b(1)
  &=33\,\varepsilon_{\rm round}^{(0)}+6\cdot2^{-256},\\
 b(2^{128})
  &=\left(1+2^{-123}\right)\varepsilon_{\rm round}^{(0)}
    +3(2^{128}+2^{-128})2^{-256}.
\end{align*}
This yields 101.74641188359571 and 106.79080421773332 bits.  Applying the
release's whole-ledger factor three subtracts \(\log_2 3\) from each endpoint,
giving
\[
 100.16144938287455\quad\text{and}\quad105.20584171701216
\]
bits.  The first is the published floor.

The selector-local containment in Lemma~\ref{lem:selector-soundness} gives a
strictly sharper diagnostic calculation by applying the factor three only to
the final query-miss row.  It is not used to raise the theorem's conservative
headline value.

\section{\texorpdfstring{\(\mathsf{Good}_{23}\)}{Good23} certificate structure}
\label{app:good23-certificate}

The complete product consists of the following three nonzero minors.
\begin{table}[ht]
\centering
\caption{\(\mathsf{Good}_{23}\) provenance anchors.}
\label{tab:good23-fingerprints}
\begin{tabular}{@{}lrl@{}}
\toprule
component & rank & fingerprint \\
\midrule
root-neutral joint minor & 1404 & \texttt{0x6b3838662fbf34db} \\
remaining \(G/D\) terminal Schur minor & 12 & \texttt{0x1f34525db611d292} \\
inactive-balanced \(H_1\) terminal Schur minor & 12 & \texttt{0xe70215f0b4795f52} \\
bound product & --- & \texttt{0xfc0706f3a304ae26} \\
\bottomrule
\end{tabular}
\end{table}

The first minor is selected from 10610 candidate \(\mathbb F\)-sources in
minimum query-degree order.  It contains 1068 degree-one query columns and
336 degree-two columns.  Its query, root-neutral \(z\), and
\(\gamma\)-coordinate degree bounds are 31320, 41040, and 80688.  Each Schur
minor contributes 120 to the \(z\)-degree, producing the complete bounds
\[
 d_q=31320,\qquad d_z=41280,\qquad
 d_\gamma=80688,\qquad d_{\rm cont}=121968.
\]

For a candidate schedule, the runtime checker performs the following steps:
\begin{enumerate}
  \item verify the static profile, layout, generator order, nonzero
        \(\gamma\), query count, query distinctness, and degree schema;
  \item reconstruct every source column and target coordinate from the public
        layout and schedule;
  \item verify pointwise root-neutrality and compressed-terminal-zero source
        guards;
  \item eliminate the raw query blocks, compute the two terminal Schur
        complements, and require the ranks in
        Table~\ref{tab:good23-ranks};
  \item construct the root-neutral initial/terminal/kernel map, recompute its
        pivots and nonzero minor, and require rank 1404; and
  \item bind the dynamic component reports into the complete degree/product
        report.
\end{enumerate}
The frozen fingerprints are reproducibility anchors for one known set of
pivots.  Runtime elimination may select different full-rank minors and does
not accept a supplied fingerprint in place of rank.  Transfer to another
proof instance depends on the common frozen layout and public schedule maps,
not equality of witness or proof bytes.

\section{Complete EPRO arithmetic}
\label{app:epro-ledger}

The field expander count is
\[
 2(22850+4\cdot1024)=53892.
\]
The transcript cap is
\[
 292=67\cdot4+3\cdot8
\]
possible squeeze inputs: 67 bounded extension-field challenges use at most
four blocks each, and each of the three q18 branches uses at most eight query
blocks.  Accounting for a matching state advance at every squeeze gives
\[
 637=47+2\cdot292+6.
\]
Thus the full per-attempt inventory is
\[
 C=3\cdot305152+53892+8+637=969993.
\]

At \(A=17\) and \(Q_H=2^{128}\),
\begin{align*}
 -\log_2\left(\frac{ACQ_H}{2^{256}}\right)
    &=104.02492234825198,\\
 -\log_2\left(\frac{\binom{AC}{2}}{2^{256}}\right)
    &=209.04984478399368,\\
 -\log_2\left(6A e^{-2^{27}}\right)
    &=193635243.91255844.
\end{align*}
The pairwise game traverses the simulator twice and therefore loses exactly
one bit under the triangle inequality.

\section{Security games at a glance}
\label{app:security-games}

\begin{table}[ht]
\centering
\small
\caption{Quantifiers and outputs of the three security statements.}
\label{tab:security-games}
\begin{tabular}{@{}
  >{\raggedright\arraybackslash}p{0.19\linewidth}
  >{\raggedright\arraybackslash}p{0.20\linewidth}
  >{\raggedright\arraybackslash}p{0.31\linewidth}
  >{\raggedright\arraybackslash}p{0.18\linewidth}@{}}
\toprule
game & quantified witnesses & output/view & bound \\
\midrule
argument soundness & no witness for chosen \(x\) & accepting proof &
100.161 bits per random-oracle query \\
real versus simulator & one valid \(w\) for fixed \(x\) & complete Proof/Abort &
104.024 bits \\
pairwise hiding & valid \(w_0,w_1\) for fixed \(x\) & complete Proof/Abort &
103.024 bits \\
\bottomrule
\end{tabular}
\end{table}

The hiding games allow arbitrary auxiliary input and adaptive post-output
queries within \(Q_H\le2^{128}\).  The soundness game is classical and uses
the same query ceiling.  The exact environmental scope of these games is
collected in Section~\ref{sec:limitations}.
